\documentclass[11pt]{article}
\usepackage[utf8]{inputenc}
\usepackage[T1]{fontenc}
\usepackage{amsmath, amssymb, amsthm}
\usepackage{graphicx}
\usepackage{xcolor}
\usepackage{hyperref}

\newtheorem{theorem}{Theorem}[section]
\newtheorem{lemma}[theorem]{Lemma}
\newcommand{\R}{\mathbb{R}}
\newcommand{\norm}[1]{\left\lVert #1 \right\rVert}

\title{Fourier Series in Practice}
\author{Ada Lovelace \and Carl Friedrich Gauss}
\date{March 2025}

\begin{document}
\maketitle

\begin{abstract}
We review how periodic signals are written as sums of sines and cosines and
show a few results with \emph{numbered equations}, a table and a figure.
\end{abstract}

\section{Introduction}
\label{sec:intro}
A function $f\colon \R \to \R$ with period $2\pi$ can be expanded as
\begin{equation}
  f(x) = \frac{a_0}{2} + \sum_{n=1}^{\infty} \bigl( a_n \cos nx + b_n \sin nx \bigr).
  \label{eq:fourier}
\end{equation}
The coefficients in~\eqref{eq:fourier} are given by
\begin{align}
  a_n &= \frac{1}{\pi} \int_{-\pi}^{\pi} f(x) \cos nx \, dx, \\
  b_n &= \frac{1}{\pi} \int_{-\pi}^{\pi} f(x) \sin nx \, dx. \nonumber
\end{align}

\subsection{Convergence}
\begin{theorem}[Dirichlet]
If $f$ is piecewise smooth, its Fourier series converges to $\frac{1}{2}\bigl(f(x^+) + f(x^-)\bigr)$.
\end{theorem}

\begin{proof}
See any textbook, for example~\cite{katznelson}.
\end{proof}

\section{Examples}
\label{sec:examples}
\begin{itemize}
  \item The square wave has only odd harmonics.
  \item The sawtooth wave has all harmonics with $b_n = \frac{2(-1)^{n+1}}{n}$.
  \item Energy: $\norm{f}^2 = \frac{a_0^2}{2} + \sum_n (a_n^2 + b_n^2)$.
\end{itemize}

\begin{table}[h]
  \centering
  \begin{tabular}{|l|c|r|}
    \hline
    Signal & Harmonics & Amplitude \\
    \hline
    Square & odd & $4/(\pi n)$ \\
    Sawtooth & all & $2/n$ \\
    \hline
    \multicolumn{3}{|c|}{Triangle: odd, $8/(\pi n)^2$} \\
    \hline
  \end{tabular}
  \caption{Harmonic content of common signals.}
  \label{tab:signals}
\end{table}

\begin{figure}[h]
  \centering
  \includegraphics[width=0.6\textwidth]{figures/square-wave.png}
  \caption{Partial sums of the square wave.}
  \label{fig:square}
\end{figure}

Table~\ref{tab:signals} and Figure~\ref{fig:square} summarize Section~\ref{sec:examples}.%
\footnote{All examples use $L = \pi$.}
The matrix
\[
  M = \begin{pmatrix} \cos\theta & -\sin\theta \\ \sin\theta & \cos\theta \end{pmatrix}
\]
rotates the plane, and \textcolor{teal}{colored text} or \href{https://en.wikipedia.org/wiki/Fourier_series}{a link} work too.

\begin{thebibliography}{9}
\bibitem{katznelson} Y.~Katznelson, \emph{An Introduction to Harmonic Analysis}, Cambridge, 2004.
\end{thebibliography}

\end{document}
